%% kapa: exam=1 round=33 year=2022 period=2 form=A subject=appraisal_law subject_ko=감정평가관계법규 seq=38 number=38
%% source: Q-Net 감정평가사 기출문제 (cst00309) artlSeq=5209997 file=raw/1008/5209997/2022년도 제33회 감정평가사 1차시험 2교시(감정평가 관계법규, 회계학).pdf
%% conversion: kordoc 4.24.51 -> tools/gen1.py (auto); review=proofread-sonnet; flags=none
\begin{question}{38}
도시 및 주거환경정비법령상 시장·군수등이 직접 정비사업을 시행하거나 토지주택공사등을 사업시행자로 지정하여 정비사업을 시행하게 할 수 있는 경우에 해당하지 \uline{않는} 것은?
\begin{choices}
  \item 천재지변으로 긴급하게 정비사업을 시행할 필요가 있다고 인정하는 때
  \item 재건축조합이 사업시행 예정일부터 2년 이내에 사업시행계획인가를 신청하지 아니한 때
  \item 조합설립추진위원회가 시장·군수등의 구성승인을 받은 날부터 3년 이내에 조합설립인가를 신청하지 아니한 때
  \item 지방자치단체의 장이 시행하는 「국토의 계획 및 이용에 관한 법률」에 따른 도시·군계획사업과 병행하여 정비사업을 시행할 필요가 있다고 인정하는 때
  \item 해당 정비구역의 국·공유지 면적 또는 국·공유지와 토지주택공사등이 소유한 토지를 합한 면적이 전체 토지면적의 2분의 1 이상으로서 토지등소유자의 과반수가 시장·군수등 또는 토지주택공사등을 사업시행자로 지정하는 것에 동의하는 때
\end{choices}
\end{question}
\answer{2}
